\part{Aktueller Stand}

\chapter{Was Modelle leisten}

\section{Eine Bilanz}

\begin{merksatz}
\boxtitel{Die vierzehn Kapitel, geordnet nach dem, was sie leisten}
\begin{center}
{\small
\renewcommand{\arraystretch}{1.5}
\begin{tabular}{@{}P{2.8cm}P{2.2cm}P{3.2cm}@{}}
\toprule
\textbf{Leistung} & \textbf{Kapitel} & \textbf{Beispiel}\\
\midrule
Größen abschätzen & $1$, $10$, $13$ & vier Halbwertszeiten\\
Schwellen finden & $7$, $8$, $12$ & $\Rnull = 1$\\
Überraschungen & $3$, $4$, $14$ & Volterras Prinzip\\
Annahmen prüfen & $5$, $9$, $11$ & Durchmischung\\
Einsicht erzeugen & $6$, $12$ & Schwelle ohne Schwellenwert\\
\bottomrule
\end{tabular}}
\end{center}

\begin{center}
\begin{minipage}{0.88\textwidth}
\kursiv{Bemerkenswert ist, was \kursiv{nicht} in der Liste steht: genaue Vorhersagen. Kein Modell dieses Bandes hat eine Zahl geliefert, auf die man sich verlassen könnte. Die Leistung liegt woanders.}
\end{minipage}
\end{center}
\end{merksatz}

\section{Vier Muster, die überall wiederkehrten}

\begin{merksatz}
\boxtitel{Muster 1 --- die Durchmischungsannahme}
\begin{center}
{\small
\renewcommand{\arraystretch}{1.48}
\begin{tabular}{@{}P{1.8cm}P{3.0cm}P{3.4cm}@{}}
\toprule
\textbf{Kapitel} & \textbf{Wo} & \textbf{Folge der Verletzung}\\
\midrule
$4$ & Räuber trifft Beute & Zyklen schwächer\\
$5$ & Arten konkurrieren & Koexistenz möglich\\
$7$ & Jeder trifft jeden & Endgröße überschätzt\\
$9$ & Kontaktnetze & Schwelle verschoben\\
$11$ & Enzym trifft Substrat & Umsatz anders\\
$14$ & Klone konkurrieren & Eindämmung schwächer\\
\bottomrule
\end{tabular}}
\end{center}

\begin{center}
\begin{minipage}{0.9\textwidth}
\textbf{Sechs Kapitel, dieselbe schwächste Stelle.} Das Massenwirkungsgesetz stammt aus der Chemie, wo es für gerührte Lösungen gilt. In der Biologie ist nichts gerührt.
\end{minipage}
\end{center}

\begin{center}
\begin{minipage}{0.88\textwidth}
\kursiv{Und die Richtung des Fehlers ist meist dieselbe: Die Annahme überschätzt die Zahl der Begegnungen zwischen \kursiv{verschiedenen} Gruppen. Ausbreitung, Konkurrenz und Umsatz fallen in der Wirklichkeit schwächer aus als berechnet.}
\end{minipage}
\end{center}
\end{merksatz}

\begin{merksatz}
\boxtitel{Muster 2 --- verzögerte Rückkopplung erzeugt Schwingungen}
\begin{center}
{\small
\renewcommand{\arraystretch}{1.48}
\begin{tabular}{@{}P{1.8cm}P{5.8cm}@{}}
\toprule
$3$ & Diskrete Zeitschritte statt stetiger Änderung --- Zyklen und Chaos.\\
$4$ & Der Räuber folgt der Beute mit einem Viertel Umlauf Verzögerung.\\
$8$ & Impfbereitschaft folgt der wahrgenommenen Gefahr, die der Krankheitslast folgt.\\
$12$ & Die Erholungsgröße $w$ folgt der Spannung langsam --- das Aktionspotential.\\
\bottomrule
\end{tabular}}
\end{center}

\begin{center}
\begin{minipage}{0.9\textwidth}
\textbf{Eine Rückkopplung, die sofort wirkt, stabilisiert. Eine, die verzögert wirkt, schwingt.} Das ist der einzige Satz, den man aus diesem Band mitnehmen sollte, wenn man nur einen mitnimmt.
\end{minipage}
\end{center}
\end{merksatz}

\begin{merksatz}
\boxtitel{Muster 3 --- schnell und langsam}
\begin{center}
{\small
\renewcommand{\arraystretch}{1.48}
\begin{tabular}{@{}P{1.8cm}P{2.6cm}P{3.6cm}@{}}
\toprule
\textbf{Kap.} & \textbf{schnell} & \textbf{langsam}\\
\midrule
$10$ & Verteilung & Ausscheidung\\
$11$ & Komplexbildung & Substratverbrauch\\
$12$ & Spannung & Erholung\\
$13$ & Zelltötung & Nachwachsen\\
\bottomrule
\end{tabular}}
\end{center}

\begin{center}
\begin{minipage}{0.88\textwidth}
\kursiv{Überall dort erlaubt die Trennung eine Vereinfachung: Der schnelle Vorgang wird als augenblicklich eingestellt behandelt --- die Quasi-Stationarität aus Kapitel~11, die Verteilungsphase aus Kapitel~10, die Nullkline aus Kapitel~12. Diese Trennung ist das mächtigste Vereinfachungswerkzeug der Biomathematik.}
\end{minipage}
\end{center}
\end{merksatz}

\begin{merksatz}
\boxtitel{Muster 4 --- Schwellen}
\begin{center}
{\small
\renewcommand{\arraystretch}{1.48}
\begin{tabular}{@{}P{1.8cm}P{2.4cm}P{3.8cm}@{}}
\toprule
\textbf{Kap.} & \textbf{Schwelle} & \textbf{Was sie trennt}\\
\midrule
$3$ & $a = 3$ & Einpendeln / Zyklus\\
$5$ & $\alpha_{12}\alpha_{21} = 1$ & Koexistenz / Verdrängung\\
$6$ & $T = 0$ & Ruhe / Schwingung\\
$7$ & $\Rnull = 1$ & Erlöschen / Epidemie\\
$12$ & $v^{*2} = 1-\varepsilon b$ & Ruhe / Dauerfeuer\\
$14$ & $\mu N = 1$ & Heilbar / nicht\\
\bottomrule
\end{tabular}}
\end{center}

\begin{center}
\begin{minipage}{0.9\textwidth}
\textbf{Alle sechs Schwellen sind Bifurkationen im Sinne von Kapitel~6.} Sie trennen nicht wenig von viel, sondern das eine Verhalten vom anderen.
\end{minipage}
\end{center}

\begin{center}
\begin{minipage}{0.88\textwidth}
\kursiv{Darin liegt die eigentliche Leistung der Modellierung: Ein Modell findet Schwellen, die man ohne es nicht sähe --- und an denen eine kleine Änderung der Bedingungen eine große Änderung des Verhaltens bewirkt. Das ist mehr wert als jede Zahlenvorhersage.}
\end{minipage}
\end{center}
\end{merksatz}

\section{Das Parameterproblem}

\begin{warnung}
\boxtitelr{Die Grenze, die immer wiederkam}
\begin{center}
{\small
\renewcommand{\arraystretch}{1.5}
\begin{tabular}{@{}P{1.8cm}P{6.2cm}@{}}
\toprule
$2$ & $K$ ist vor dem Wendepunkt nicht schätzbar --- Pearls Fehlvorhersage.\\
$7$ & $\pm20\,\%$ in $\Rnull$ verdoppeln die Gipfelhöhe.\\
$9$ & $\langle k^2\rangle$ hängt an den wenigen extremen Werten.\\
$11$ & $V_{\max}$ wird nie erreicht und daher schlecht geschätzt.\\
$13$ & Halbierte Verdopplungszeit kehrt die Antwort um.\\
\bottomrule
\end{tabular}}
\end{center}

\begin{center}
\begin{minipage}{0.9\textwidth}
\textbf{Die Gleichungen sind nicht das Problem. Die Zahlen darin sind es.}
\end{minipage}
\end{center}

\schritt{Die drei Arten von Schwierigkeit}
\begin{center}
{\small
\renewcommand{\arraystretch}{1.5}
\begin{tabular}{@{}P{3.0cm}P{5.6cm}@{}}
\toprule
\textbf{Nicht messbar} & Die Größe ist im Modell wohldefiniert und in der Wirklichkeit nicht auffindbar --- etwa die Kapazitätsgrenze eines Tumors.\\
\textbf{Nicht bestimmbar} & Verschiedene Parameterkombinationen liefern dieselben Daten. Man nennt das fehlende \Begriff{Identifizierbarkeit}.\\
\textbf{Nicht konstant} & Die Größe ändert sich mit Bedin\-gungen, Zeit und Verhalten --- die erste Warnung aus dem Vorwort.\\
\bottomrule
\end{tabular}}
\end{center}

\begin{center}
\begin{minipage}{0.88\textwidth}
\kursiv{Die mittlere Zeile lässt sich prüfen, bevor man Daten erhebt: Man untersucht, ob die Abbildung von den Parametern auf die beobachtbaren Größen umkehrbar ist. Ein Modell, das diese Prüfung nicht besteht, liefert Parameterschätzungen mit engen Vertrauensbereichen, die nichts bedeuten --- ein Fehler, der in der Literatur häufig ist.}
\end{minipage}
\end{center}
\end{warnung}

\begin{warnung}
\boxtitelr{Und der Preis der Verfeinerung}
Jede zusätzliche Kammer, jede Altersgruppe, jeder Mechanismus bringt neue Parameter.

\begin{center}
\renewcommand{\arraystretch}{1.5}
\begin{tabular}{@{}P{2.6cm}P{2.4cm}P{2.6cm}@{}}
\toprule
\textbf{Modell} & \textbf{Parameter} & \textbf{Anpassung}\\
\midrule
SIR & $2$ & \textcolor{mittel}{grob}\\
SEIR mit Altersgruppen & $\sim 20$ & besser\\
Vollständiges Netzwerkmodell & $> 100$ & \textcolor{gruen}{sehr gut}\\
\bottomrule
\end{tabular}
\end{center}

\begin{center}
\begin{minipage}{0.9\textwidth}
\textbf{Die Anpassung wird besser, die Vorhersage nicht.} Ein Modell mit hundert Parametern lässt sich an fast jede Beobachtung anpassen --- und sagt gerade deshalb wenig.
\end{minipage}
\end{center}

\begin{center}
\begin{minipage}{0.88\textwidth}
\kursiv{Das ist die dritte Warnung aus dem Vorwort, bestätigt. Sie ist keine Aufforderung zur Einfachheit um jeden Preis: Manche Fragen --- etwa die Wirkung von Schulschließungen --- lassen sich ohne Altersstruktur gar nicht stellen. Die Regel lautet: so viel Verfeinerung wie die Frage verlangt, und keine mehr.}
\end{minipage}
\end{center}
\end{warnung}

\section{Wenn das Modell die Wirklichkeit ändert}

\begin{merksatz}
\boxtitel{Eine Eigenart, die es in der Physik nicht gibt}
\begin{center}
\begin{minipage}{0.9\textwidth}
\textbf{Ein Planet ändert seine Bahn nicht, weil jemand sie berechnet hat. Eine Bevölkerung ändert ihr Verhalten, wenn eine Vorhersage veröffentlicht wird.}
\end{minipage}
\end{center}

\schritt{Die Folge}
\begin{center}
{\small
\renewcommand{\arraystretch}{1.5}
\begin{tabular}{@{}P{2.8cm}P{5.8cm}@{}}
\toprule
\textbf{Wirkt die Warnung} & Die Vorhersage trifft nicht ein --- \kursiv{weil} sie gemacht wurde.\\
\textbf{Wird sie ignoriert} & Sie trifft ein, und man wirft der Modellierung vor, zu spät gewarnt zu haben.\\
\textbf{Beurteilung} & In beiden Fällen lässt sich nicht mehr feststellen, ob das Modell gut war.\\
\bottomrule
\end{tabular}}
\end{center}

\begin{center}
\begin{minipage}{0.88\textwidth}
\kursiv{Das ist kein rhetorisches Problem, sondern ein logisches: Ein Modell, dessen Vorhersage das Verhalten ändert, sagt nicht vorher, was geschehen wird, sondern was geschähe, wenn nichts geschähe. Das ist eine andere Aussage --- und sie wird regelmäßig als die erste gelesen.}
\end{minipage}
\end{center}

\schritt{Der Umgang damit}
\begin{center}
\renewcommand{\arraystretch}{1.5}
\begin{tabular}{@{}P{2.4cm}P{6.2cm}@{}}
\toprule
\textbf{Szenarien} & Statt einer Vorhersage rechnet man mehrere Verläufe unter verschiedenen Annahmen über das Verhalten.\\
\textbf{Bänder} & Unsicherheitsbereiche statt einzelner Kurven.\\
\textbf{Benennen} & Ausdrücklich sagen, welche Annahmen über Verhalten eingegangen sind.\\
\bottomrule
\end{tabular}
\end{center}

\begin{center}
\begin{minipage}{0.88\textwidth}
\kursiv{Die Ensemble-Rechnung aus Kapitel~3 ist der erprobte Umgang mit unsicheren Anfangswerten. Für unsicheres \kursiv{Verhalten} gibt es kein entsprechend ausgereiftes Verfahren --- und das ist eine der offenen Fragen des Gebiets.}
\end{minipage}
\end{center}
\end{merksatz}

\section{Mechanismus oder Vorhersage}

\begin{merksatz}
\boxtitel{Zwei Zugänge}
\begin{center}
{\small
\renewcommand{\arraystretch}{1.52}
\begin{tabular}{@{}P{2.0cm}P{2.9cm}P{3.3cm}@{}}
\toprule
& \textbf{Mechanistisch} & \textbf{Daten\-getrieben}\\
\midrule
Aufbau & aus Annahmen & aus Beispielen\\
Parameter & wenige, deutbar & viele, nicht deutbar\\
Übertragbar & \textcolor{gruen}{auf neue Lagen} & \textcolor{akzent}{nur auf ähnliche}\\
Eingriffe & \textcolor{gruen}{vorhersagbar} & \textcolor{akzent}{nicht}\\
Genauigkeit & \textcolor{mittel}{mäßig} & \textcolor{gruen}{oft hoch}\\
\bottomrule
\end{tabular}}
\end{center}

\begin{center}
\begin{minipage}{0.88\textwidth}
\kursiv{Die vierte Zeile ist die entscheidende. Ein Modell, das nur aus Daten gelernt hat, kann nicht sagen, was geschähe, wenn man eingreift --- weil es keine Beispiele für den Eingriff gesehen hat. Genau das aber will man in der Medizin wissen. Das ist dieselbe Unterscheidung wie zwischen Vorhersage und Kausalität in Band 17, Kapitel~15.}
\end{minipage}
\end{center}

\schritt{Der Mittelweg}
Man baut die bekannte Struktur ein und lässt nur das Unbekannte lernen. Der Erhaltungssatz aus Kapitel~7, die Sättigungsform aus Kapitel~11, die Kammerstruktur aus Kapitel~10 werden vorgegeben; was innerhalb dieser Struktur unbekannt ist, wird aus Daten geschätzt.

\begin{center}
\begin{minipage}{0.88\textwidth}
\kursiv{Solche Verfahren sind seit etwa 2019 ein aktives Feld. Ihr Vorteil ist, dass sie mit weniger Daten auskommen und deutbare Größen liefern. Ihre Schwierigkeit: Die vorgegebene Struktur muss stimmen --- sonst lernt das Verfahren die Abweichung als Rauschen.}
\end{minipage}
\end{center}
\end{merksatz}

\section{Wo das Fach heute steht}

\begin{merksatz}
\boxtitel{Die offenen Fragen}
\begin{center}
{\small
\renewcommand{\arraystretch}{1.5}
\begin{tabular}{@{}P{3.2cm}P{5.4cm}@{}}
\toprule
\textbf{Verhalten einbeziehen} & Modelle, in denen Menschen auf die Lage reagieren --- unausgereift.\\
\textbf{Vielskalig} & Molekül, Zelle, Organ, Mensch, Bevölkerung in einem Modell --- rechnerisch möglich, methodisch offen.\\
\textbf{Unsicherheit berichten} & Wie berichtet man Bän\-der so, dass sie verstanden werden?\\
\textbf{Modellwahl} & Wann ist ein Modell gut genug? Es gibt keine allgemeine Antwort.\\
\bottomrule
\end{tabular}}
\end{center}

\begin{center}
\begin{minipage}{0.88\textwidth}
\kursiv{Die dritte Zeile ist weniger mathematisch, als sie klingt, und vielleicht die wichtigste. Ein Band, das als Bereich berichtet wird, wird in der Öffentlichkeit regelmäßig auf seinen ungünstigsten Rand verkürzt --- oder auf seine Mitte. Beides verfälscht.}
\end{minipage}
\end{center}
\end{merksatz}

\begin{merksatz}
\boxtitel{Sechs Sätze, die bleiben}
\begin{center}
{\small
\renewcommand{\arraystretch}{1.55}
\begin{tabular}{@{}P{0.7cm}P{7.7cm}@{}}
\toprule
$1$ & Nichts wächst lange exponentiell. Wer eine Wachstumsrate hochrechnet, muss sagen, wo das Modell endet.\\
$2$ & Verzögerte Rückkopplung erzeugt Schwingungen. Sofortige stabilisiert.\\
$3$ & Eine Schwelle entsteht aus der Form der Gleichungen, nicht aus einem eingebauten Schwellenwert.\\
$4$ & Nicht der Mittelwert bestimmt die Ausbreitung, sondern die Streuung.\\
$5$ & Ein Modell, das man an alles anpassen kann, sagt nichts vorher.\\
$6$ & Jedes Modell ist falsch. Die brauchbaren sind die, deren Grenzen man kennt.\\
\bottomrule
\end{tabular}}
\end{center}
\end{merksatz}

\begin{merksatz}
\boxtitel{Ein Schlusswort}
\begin{center}
\begin{minipage}{0.9\textwidth}
Dieser Band begann mit der Feststellung, dass jedes seiner Modelle falsch ist, und mit der Frage, welche Auskunft sie trotzdem geben.

Fünfzehn Kapitel später lässt sich die Antwort benennen. Kein Modell dieses Bandes hat vorhergesagt, was geschehen wird. Jedes hat gezeigt, \kursiv{woran es liegt} --- warum eine Kultur sättigt, warum ein Eingriff der Beute nützt, warum eine Epidemie über ihre Schwelle hinausschießt, warum eine Nervenzelle schaltet, obwohl nichts in ihr schaltet, warum die Behandlung mit der besseren Anfangswirkung die schlechtere sein kann.

Das ist weniger, als man von einem Modell erhofft, und mehr, als man ohne Mathematik bekäme. Die Einsichten dieses Bandes sind nicht in Zahlen formulierbar, sondern in Bedingungen: \kursiv{wenn} die Rückkopplung verzögert ist, \kursiv{dann} schwingt es. \kursiv{Wenn} die Streuung groß ist, \kursiv{dann} sinkt die Schwelle. Solche Sätze gelten unabhängig davon, wie gut die Parameter bekannt sind.
\end{minipage}
\end{center}

\begin{center}
\begin{minipage}{0.88\textwidth}
\kursiv{Und sie sind der Grund, warum die Biomathematik trotz ihrer unsicheren Zahlen ein ernstes Fach ist. Sie liefert nicht die Antwort, sondern die Frage in der Form, in der sie beantwortbar wird --- und weist dabei regelmäßig auf Zusammenhänge hin, die ohne sie unsichtbar blieben. Das ist die Leistung, um die es von Anfang an ging.}
\end{minipage}
\end{center}
\end{merksatz}
